% ID: MAT-TRI-EQG-007
% Titolo: Valutare affermazioni goniometriche (2)
% Disciplina: Matematica
% Area: trigonometria
% Argomento: equazioni_goniometriche
% Classe: Quarta secondaria di secondo grado
% Difficolta: 4
% Tipo: vero_falso
% Tag: matematica, trigonometria, equazioni_goniometriche
% Fonte: verifica_goniometria_intermedia_anonimizzata
% Provenienza: materiale fornito dal docente; testo anonimizzato; soluzione redatta per Exergo

\begin{esercizio}
Vero o falso? Giustifica brevemente. 

\begin{enumerate}
\item L'equazione $\sin 3x=\sin x$ ammette esattamente sei
soluzioni nell'intervallo $[0;2\pi)$. 

\item Ogni soluzione di $\sin 2x=0$ è anche soluzione di
$\sin x=0$. 

\item L'uguaglianza $\dfrac{\sin 2x}{\sin x}=2\cos x$
è verificata per ogni $x\in\mathbb{R}$. 

\item L'equazione $1-\cos 2x=\sin x$ ha quattro soluzioni
in $[0;2\pi)$. 
\end{enumerate}
Salvo diversa indicazione, lavora in $\mathbb R$; angoli in radianti.
\end{esercizio}

\begin{soluzione}
a) Vero: da $\sin3x=3\sin x-4\sin^3x$ si ottiene $2\sin x(1-2\sin^2x)=0$.
In $[0,2\pi)$ le soluzioni sono $0,\pi,\pi/4,3\pi/4,5\pi/4,7\pi/4$.
b) Falso: $x=\pi/2$ è un controesempio.
c) Falso: occorre escludere $x=k\pi$.
d) Vero: $2\sin^2x=\sin x$ dà $0,\pi,\pi/6,5\pi/6$.
\end{soluzione}
